\chapter{Principe d'Ambivalence : orientations conjuguées \texorpdfstring{\(\pi/\varphi^2\)}{π/φ²} et \texorpdfstring{\(\varphi/\pi^2\)}{φ/π²}, involution des feuillets et réalisation euclidienne du principe d'équivalence.}
\label{chap:83-16}

\section{Sections compatibles et théorie de la mesure}
\label{p12-83-c16-s1:chap:medida}

La notion dimensionnelle de mesure se formule sur des systèmes inverses d'états
compatibles. Avant d'attribuer une valeur chiffrée à un état, on spécifie son histoire
de raffinement, le couplage avec l'appareil et l'application qui produit
l'observable. Cette séquence permet de séparer l'existence d'une valeur réelle de
l'opération physique par laquelle cette valeur est enregistrée.

\subsection{Le nombre comme section compatible}

Soit \((\mathcal S_n,r_{n+1,n})_{n\geq0}\) un système inverse d'états
survivants. Une section compatible est une suite
\(x=(x_n)_{n\geq0}\) telle que
\[
 r_{n+1,n}(x_{n+1})=x_n
 \qquad(n\geq0).
\]
Son espace est la limite inverse
\[
 \mathcal X_\infty=\varprojlim_n\mathcal S_n.
\]
Lorsque chaque \(x_n\) détermine un intervalle \(I_n(x_n)\subset\mathbb R\), que les
adhérences \(\overline{I_n(x_n)}\) sont compactes non vides, emboîtées et que
leurs diamètres tendent vers zéro, le théorème des intervalles emboîtés garantit
que leur intersection contient un point unique. L'évaluation archimédienne est
alors définie par
\[
 \mathcal P_{\rm arq}(x)
 =\text{l'unique élément de }
   \bigcap_{n\geq0}\overline{I_n(x_n)}.
\]
Le nombre HMT est la section compatible
\(x\in\mathcal X_\infty\). Le lecteur \(L\) est une structure supplémentaire, et le
couplage \((x,L)\) constitue sa présentation lue ou mesure. Deux
sections peuvent posséder la même évaluation réelle et différer dans leurs coordonnées
de quotient, d'orientation ou de frontière ; la projection archimédienne n'identifie pas
par définition ces généalogies.

\subsection{Compatibilité entre cylindres ternaires et millièmes}

Un mot ternaire de longueur \(L\), dont la valeur entière est \(N\), détermine
le cylindre
\[
 I_{N,L}=\left[\frac{N}{3^L},\frac{N+1}{3^L}\right).
\]
Un mot de \(m\) blocs en base \(1000\), dont la valeur entière est \(D\),
détermine
\[
 J_{D,m}=\left[\frac{D}{1000^m},\frac{D+1}{1000^m}\right).
\]

\begin{lemma}[Compatibilité exacte des cylindres]
Les cylindres \(I_{N,L}\) et \(J_{D,m}\) ont une intersection non vide si et
seulement si
\[
 N\,1000^m<(D+1)3^L,
 \qquad
 D\,3^L<(N+1)1000^m.
\]
\end{lemma}
\begin{proof}
Deux intervalles semi-ouverts \([a,b)\) et \([c,d)\) se coupent exactement
lorsque \(a<d\) et \(c<b\). En substituant les quatre extrémités et en multipliant
par l'entier positif \(3^L1000^m\), on obtient les inégalités. La
décision se réduit donc à l'arithmétique entière.
\end{proof}

Le prolongement d'un mot raffine son cylindre. Une chaîne compatible dont les
adhérences sont compactes, non vides et emboîtées, avec un diamètre convergeant vers zéro,
détermine une unique évaluation réelle. L'emploi des adhérences est essentiel :
les intervalles semi-ouverts emboîtés peuvent avoir une intersection vide lorsque la
limite coïncide avec une extrémité exclue. L'état complet conserve les
mots, les transformations et les retenues qui produisent l'évaluation.

\subsection{Couplage réversible entre système et appareil}

Soient \(X\) l'espace des états du système, \(M\) l'espace des états de
l'appareil et \(m_0\in M\) sa préparation initiale. Un couplage est une
application
\[
 U:X\times M\longrightarrow X\times M.
\]
Lorsque \(U\) est bijective, l'état conjoint conserve l'information. Une
application d'évaluation \(q_M:M\to Y\) peut être non injective et ne publier
qu'un observable \(y\in Y\).

\begin{definition}[Mesure]
Une mesure HMT est le triplet
\[
 \mathfrak M=(U,q_M,\mathcal C),
\]
où \(U\) couple système et appareil, \(q_M\) évalue l'état de l'appareil et
\(\mathcal C\subseteq X\times M\) est la condition de compatibilité. Sa
sortie sur \(x\) est
\[
 \operatorname{Med}_{\mathfrak M}(x)
 =q_M\!\left(\operatorname{pr}_M U(x,m_0)\right),
\]
pourvu que le couple couplé satisfasse \(\mathcal C\).
\end{definition}

La définition établit une composition précise de trois opérations :
couplage, restriction au domaine compatible et évaluation de l'observable.

\begin{theorem}[Localisation de la perte d'information]
Soit
\[
 \Omega_M:=q_M\circ\operatorname{pr}_M:X\times M\longrightarrow Y.
\]
Si \(U\) est bijective et si l'état final complet est conservé, les classes
d'entrées indiscernables par l'observable publié sont exactement les
images réciproques par \(U\) des fibres de \(\Omega_M\), sauf si l'on compose
en outre avec une suppression ou une réinitialisation ultérieure de l'état.
\end{theorem}
\begin{proof}
La bijectivité de \(U\) permet de récupérer \((x,m_0)\) à partir de l'état final
conjoint. Pour deux entrées \(z,z'\) du domaine compatible, l'égalité des
registres équivaut à
\(\Omega_M(Uz)=\Omega_M(Uz')\). Par conséquent, \(Uz\) et \(Uz'\) appartiennent à une
même fibre de l'application composée \(q_M\circ\operatorname{pr}_M\). Cette
composition inclut aussi bien la perte due à la projection sur l'appareil que
celle due à l'évaluation de son état. Tout effacement ultérieur introduit une
application supplémentaire et doit être analysé séparément.
\end{proof}

\begin{figure}[htbp]
\centering
\resizebox{\textwidth}{!}{%
\begin{tikzpicture}[node distance=13mm and 15mm,>=Latex,
box/.style={draw=hmtblue,rounded corners,fill=hmtlight,align=center,minimum height=9mm,minimum width=29mm}]
\node[box] (s) {état complet\\[-1mm]$x$};
\node[box,right=of s] (u) {acoplamiento\\[-1mm]$U(x,m_0)$};
\node[box,right=of u] (c) {restricción\\[-1mm]$\mathcal C$};
\node[box,right=of c] (q) {evaluación\\[-1mm]$q_M$};
\node[box,right=of q] (y) {observable\\[-1mm]$y$};
\draw[->,thick] (s)--(u); \draw[->,thick] (u)--(c);
\draw[->,thick] (c)--(q); \draw[->,thick] (q)--(y);
\end{tikzpicture}
}
\caption{Composición de una medición. La aplicación de evaluación actúa
después del acoplamiento y de la restricción de compatibilidad; el estado
conjunto permanece disponible para la transformación inversa.}
\label{p12-83-c16-s1:fig:medicion}
\end{figure}

\subsection{Principio de Ambivalencia: involución bidireccional de los lectores π/φ² y φ/π²}

Sean
\[
 q_A:\mathcal X_{\rm TPK}\to\mathcal X_A,
 \qquad
 q_V:\mathcal X_{\rm TPK}\to\mathcal X_V
\]
deux applications quotient associées, respectivement, aux observables
d'ouverture et de rétention, et soit
\(T:\mathcal X_{\rm TPK}\to\mathcal M\) la variable de mémoire
transportée.

\begin{principle}[Ambivalence]
L'état complet précède ses projections observables. Les relations
d'équivalence induites par \(q_A\) et \(q_V\) ne sont pas supposées identiques, et la
compatibilité entre les deux est enregistrée par \(T\). L'application conjointe
\[
 (q_A,q_V,T):\mathcal X_{\rm TPK}
 \longrightarrow\mathcal X_A\times\mathcal X_V\times\mathcal M
\]
est adoptée comme condition de séparation des états physiquement
discernables, modulo le groupe de jauge déclaré.
\end{principle}

\begin{proposition}[Reconstruction sous la condition de séparation]
Si l'application conjointe précédente est injective modulo jauge, chaque classe
d'état complet est déterminée de manière unique par ses deux observables et sa
variable de mémoire.
\end{proposition}
\begin{proof}
Deux états ayant le même triplet d'images appartiennent à une même classe de
jauge par l'hypothèse d'injectivité. Chaque fibre de l'application conjointe
contient donc au plus une classe.
\end{proof}

L'observable \(q_A\) enregistre les prolongements accessibles et \(q_V\) la
composante retenue. Une interprétation matérielle de la rétention exige un
observable de masse qui se factorise par \(q_V\) ; elle ne se déduit pas uniquement du
principe de séparation.

\subsection{Détermination relationnelle et orientation temporelle}

Soit \([\gamma]_{\mathcal G}\) la classe d'une trajectoire dans le groupoïde des
états et soit \(\beta_\gamma\) sa donnée terminale de frontière.
\begin{principle}[Détermination relationnelle]
Un observable local admissible a la forme
\[
 \mathcal O_{\rm loc}(x;\gamma)
 =F\bigl(x,[\gamma]_{\mathcal G},\beta_\gamma\bigr)
\]
et est invariant sous les équivalences de trajectoire qui préservent le
registre ordonné. Par conséquent, la coordonnée ponctuelle visible ne détermine pas
à elle seule l'observable lorsque l'évaluation a identifié des classes de
transport différentes.
\end{principle}

La formulation est relationnelle en un sens précis : l'argument de
\(\mathcal O_{\rm loc}\) contient l'état, la classe globale de transport et
la frontière. C'est le domaine sur lequel on peut ensuite comparer
la théorie à la tradition associée au principe de Mach. On ne présuppose pas que
toute propriété inertielle ait été dérivée ; on spécifie le type de l'application qui
doit la réaliser.

Si la mise à jour sur \(\mathcal X_{\rm TPK}\) est bijective et si le registre
conserve les coordonnées requises par son inverse, la dynamique complète est
réversible. Une orientation temporelle observable apparaît lorsqu'on la compose avec un
lecteur non injectif, qu'on réinitialise une donnée de frontière ou qu'on sélectionne un seul
prolongement parmi plusieurs. L'irréversibilité se situe donc dans un
morphisme défini de perte ou de sélection d'information, non dans le simple
ordonnancement des étiquettes.

\subsection{Relation avec les principes variationnels classiques}

Le principe de Hamilton exige que la première variation d'une action s'annule
sur la trajectoire physique ; il n'affirme pas, en général, que cette trajectoire soit un
minimum global. La fonctionnelle discrète d'interaction construite dans le
\cref{chap:pantalla}, en revanche, ordonne une famille finie de chemins selon le
nombre de changements d'orientation et de loi arithmétique. Les deux principes
ne coïncident qu'après la construction d'une limite dans laquelle le coût discret
induit une action et sa première variation reproduit les équations du
mouvement. Cette distinction empêche de remplacer une correspondance future par
une identité de noms.

\subsection{Conditions d'une évaluation prospective}

Une évaluation prospective est mathématiquement déterminée lorsque, avant de
consulter sa sortie, on fixe l'espace des états, le couplage, les
applications quotient, la variable de mémoire, la règle de compatibilité et le
critère de confrontation. Si plusieurs sections survivent, la sortie est la famille
complète. L'attribution de fréquences requiert alors une mesure de
probabilité ou une dynamique supplémentaire explicite.


\subsection{Quotient des modes et correspondance aréale--volumétrique}
\label{p12-83-c16-s2:sec:cociente-modos-fav}

Le signe du calendrier, l'étiquette de mode et le feuillet géométrique sont trois
objets typés. En particulier, on n'identifie pas le signe de phase à
l'observable angulaire sur les chiffres. Pour la correspondance géométrique, nous fixons le
lecteur HMT typé
\[
 \mathfrak g_{\rm av}(\Sigma)=\mathscr H_{\rm ar},
 \qquad
 \mathfrak g_{\rm av}(\Pi)=\mathscr H_{\rm vol}.
\]
La première attribution exprime la lecture additive ou de bord ; la seconde, la
lecture multiplicative ou d'intérieur. L'application \(\mathfrak g_{\rm av}\) fait
partie du typage des feuillets. Le résultat suivant démontre qu'une fois
ce typage fixé, la dynamique TPK impose l'incidence entre eux.

\begin{theorem}[Descente nécessaire de l'incidence des feuillets]
\label{p12-83-c16-s2:thm:descenso-necesario-fav}
Soit \(m_r\) le mode obtenu après application de la phase \(r\) du bloc
primitif \((+1,-1,0)\), avec la règle
\[
 +1:\ m\mapsto\Sigma,\qquad
 -1:\ m\mapsto\Pi,\qquad
 0:\ m\mapsto m.
\]
L'orbite cyclique des modes est \((\Sigma,\Pi,\Pi)\). Si
\[
 N_3(i,j)
 :=
 \#\{r\in\mathbb Z/3\mathbb Z:(m_r,m_{r+1})=(i,j)\},
\]
alors, dans l'ordre \((\Sigma,\Pi)\),
\[
 \boxed{
 N_3=
 \begin{pmatrix}0&1\\1&1\end{pmatrix}.}
\]
Par répétition du calendrier,
\[
 N_9=3N_3,\qquad N_{27}=9N_3,\qquad N_{108}=36N_3.
\]
Après application de \(\mathfrak g_{\rm av}\), la matrice primitive
d'incidence est, dans l'ordre
\((\mathscr H_{\rm ar},\mathscr H_{\rm vol})\),
\[
 \boxed{
 F_{\rm av}=
 \begin{pmatrix}0&1\\1&1\end{pmatrix}.}
\]
Ainsi, les deux transitions croisées, l'unique autorétention volumétrique,
l'absence d'autorétention aréale et les multiplicités primitives unitaires
sont des conséquences du quotient des modes de TPK ; ce ne sont pas quatre postulats
ajoutés.
\end{theorem}

\begin{proof}
Les trois couples consécutifs, avec fermeture cyclique, sont
\[
 \Sigma\longrightarrow\Pi,\qquad
 \Pi\longrightarrow\Pi,\qquad
 \Pi\longrightarrow\Sigma.
\]
Chacun apparaît une fois et le couple
\(\Sigma\to\Sigma\) n'apparaît pas. Cela détermine les quatre entrées de
\(N_3\). Les calendriers de longueurs \(9,27,108\) contiennent,
respectivement, \(3,9,36\) copies du bloc primitif, de sorte que leurs
matrices d'incidence sont les multiples indiqués. Le plus grand commun diviseur
des entrées non nulles de chaque matrice longue élimine uniquement cette
répétition globale et restitue \(N_3\). Enfin,
\(\mathfrak g_{\rm av}\) transporte \(\Sigma,\Pi\) vers les feuillets aréal et
volumétrique sans modifier l'incidence.
\end{proof}

La réalisation hyperbolique réelle qui utilise cette matrice et prouve son
invariance lorentzienne est formulée dans
\cref{sec:tres-generadores-concretos-trit} ; on y utilise \(F_{\rm av}\)
comme antécédent, sans répéter sa dérivation.

Ce théorème doit être distingué d'une affirmation de Markov plus forte.
L'oubli de la phase n'est pas un quotient de Markov fort de
\(\mathcal K_{\rm ph}\) : le mode visible ne détermine pas lequel des trois
opérateurs \(I,E_+,E_\times\) agit à l'itération élémentaire suivante. Ce qui descend
sans hypothèse supplémentaire est le multigraphe des arêtes d'un bloc TPK\@.

Il existe cependant une complétion à mémoire un déterminée de manière
unique par cette descente. Soit \(X_{\rm av}\) le plus petit sous-décalage de type fini
sur
\(\{\mathscr H_{\rm ar},\mathscr H_{\rm vol}\}\) qui contient tous les
mots obtenus en oubliant la phase et qui ne conserve que des couples consécutifs.
Tout sous-décalage à un pas ayant cette propriété doit admettre exactement les trois
couples observés ; le plus petit d'entre eux exclut l'unique couple non observé. Par
conséquent, sa matrice d'adjacence est \(F_{\rm av}\).

\begin{corollary}[Mesure de Parry de l'enveloppe à mémoire un]
\label{p12-83-c16-s2:cor:parry-envolvente-fav}
Le sous-décalage \(X_{\rm av}\) est primitif, a pour entropie topologique
\(\log\varphi\) et possède une unique mesure d'entropie maximale. Ses poids
stationnaires de sommet satisfont
\[
 \boxed{
 \frac{\mu_{\rm ar}}{\mu_{\rm vol}}=\varphi^{-2}.}
\]
\end{corollary}

\begin{proof}
La racine de Perron de \(F_{\rm av}\) est \(\varphi\), et un vecteur de Perron
positif est \((1,\varphi)\). Comme la matrice est symétrique, les poids de sommet
de Parry sont proportionnels au produit des vecteurs gauche et droit,
c'est-à-dire à \((1,\varphi^2)\).
\end{proof}

La mesure de Parry appartient à l'enveloppe des chemins à mémoire un. Elle n'est pas
l'image de \(\tau_{468}\), ni la fréquence temporelle de l'orbite périodique
des modes. Cette dernière est
\[
 \nu_{\rm cyc}(\mathscr H_{\rm ar})=\frac13,\qquad
 \nu_{\rm cyc}(\mathscr H_{\rm vol})=\frac23,
\]
tandis que \(\tau_{468}\) est uniforme sur les émissions et que la mesure de Parry
pondère des histoires de renouvellement de longueur arbitraire. Les trois mesures
répondent à des questions distinctes. Le rapport irrationnel
\(\varphi^{-2}\) est ainsi déduit canoniquement de l'enveloppe des chemins,
sans l'attribuer à un quotient de cardinaux finis.

\subsubsection{Du comptage chronologique au retour marqué par \texorpdfstring{\(\pi\)}{π}}
\label{p12-83-c16-s2:sec:retorno-marcado-pi-phi2}

Les deux formules
\[
 \left(\frac13,\frac23\right)
 \qquad\text{et}\qquad
 \frac{\mu_{\rm ar}}{\mu_{\rm vol}}=\varphi^{-2}
\]
ne sont pas des approximations l'une de l'autre. La première compte les trois instants du
calendrier primitif \((\Sigma,\Pi,\Pi)\) ; la seconde pondère toutes les
histoires admissibles de l'enveloppe à mémoire un. C'est pourquoi la fréquence
chronologique est rationnelle et le rapport de Parry est irrationnel. Dans ce qui suit,
nous écrivons
\[
 \frac{\pi_{\rm HMT}}{\varphi^2}
 =\pi_{\rm HMT}\varphi^{-2};
\]
il ne s'agit pas de diviser par \(\varphi^{-2}\).

L'apparition de \(\varphi^{-2}\) se voit directement dans la dynamique
linéaire. Si \(F_n\) désigne le \(n\)-ième nombre de Fibonacci,
\[
 F_{\rm av}^{\,n}
 =
 \begin{pmatrix}
  F_{n-1}&F_n\\
  F_n&F_{n+1}
 \end{pmatrix}
 \qquad(n\geq1).
\]
Les valeurs propres de \(F_{\rm av}\) sont
\(\varphi\) et \(-\varphi^{-1}\). La branche contractée inverse l'orientation en
une itération. Lors du passage à la mémoire quadratique, c'est-à-dire à
\(\operatorname{Sym}^2(\mathbb R^2)\), cette inversion est enregistrée avant la
disparition du signe. Dans la base \((x^2,2xy,y^2)\), en représentant par
lignes les images des vecteurs de base, l'opérateur induit est
\[
 \operatorname{Sym}^2(F_{\rm av})=
 \begin{pmatrix}
  0&0&1\\
  0&1&2\\
  1&1&1
 \end{pmatrix},
 \qquad
 \chi(t)=(t+1)(t^2-3t+1).
\]
Ses trois valeurs propres sont
\[
 \varphi^2,\qquad -1,\qquad \varphi^{-2}.
\]
Ainsi, \(\varphi^{-2}\) est le mode stable positif de la mémoire du second
ordre. Cette étape explique pourquoi le nombre d'or apparaît ici comme loi de
retour et non comme moyenne statistique ajoutée.

La clôture \(\pi\) appartient à un autre niveau du même état HMT : c'est la sortie
coinductive du lecteur régional déjà construit, non une nouvelle valeur propre de
\(F_{\rm av}\). Soient
\[
 P_{\rm ar}=
 \begin{pmatrix}1&0\\0&0\end{pmatrix},
 \qquad
 P_{\rm vol}=
 \begin{pmatrix}0&0\\0&1\end{pmatrix}.
\]
Parmi les opérateurs positifs qui préservent les deux feuillets, il en existe un et un
seul dont la normalisation volumétrique vaut un et dont la marque aréale est la clôture
engendrée \(\pi_{\rm HMT}\) :
\[
 D_\pi=\pi_{\rm HMT}P_{\rm ar}+P_{\rm vol}.
\]
L'unicité est relative à ces trois clauses explicites : positivité,
diagonalité par rapport aux feuillets et les deux normalisations. Avec
\(e_{\rm ar}=(1,0)^T\) et \(e_{\rm vol}=(0,1)^T\), le retour marqué de
longueur \(n\) est
\[
 \Theta_n
 :=
 \frac{\langle e_{\rm ar},
 D_\pi F_{\rm av}^{\,n}e_{\rm ar}\rangle}
 {\langle e_{\rm vol},
 F_{\rm av}^{\,n}e_{\rm vol}\rangle}
 =
 \pi_{\rm HMT}\frac{F_{n-1}}{F_{n+1}}.
\]
Par conséquent,
\[
 \boxed{\displaystyle
 \lim_{n\to\infty}\Theta_n
 =\frac{\pi_{\rm HMT}}{\varphi^2}.}
\]
La marque \(\pi_{\rm HMT}\) est appliquée une seule fois, au retour de frontière ;
elle n'est pas réinjectée à chaque transition. La combinatoire finie produit le mode
stable \(\varphi^{-2}\), le lecteur coinductif produit la clôture
\(\pi_{\rm HMT}\), et \(D_\pi\) compose les deux données sans inverser leur ordre
causal.

La longueur native \(108\) offre un témoin exact de cette convergence :
\[
 (F_{\rm av}^{108})_{\rm ar,ar}
 =F_{107}=10284720757613717413913,
\]
\[
 (F_{\rm av}^{108})_{\rm vol,vol}
 =F_{109}=26925748508234281076009.
\]
En conséquence,
\[
 \left|
 \frac{F_{107}}{F_{109}}-\varphi^{-2}
 \right|
 =
 6.1684979916614407936\ldots\times10^{-46},
\]
et, après l'unique marque de clôture,
\[
 \left|
 \pi_{\rm HMT}\frac{F_{107}}{F_{109}}
 -\frac{\pi_{\rm HMT}}{\varphi^2}
 \right|
 =
 1.9378907974286976068\ldots\times10^{-45}.
\]
Ces erreurs mesurent l'approximation finie du retour ; ce ne sont pas des ajustements de
\(\pi\) ni de \(\varphi\).

\paragraph{Lecteur cylindrique du rapport.}
Soient \(I_n^\pi=[\underline\pi_n,\overline\pi_n]\) et
\(I_n^\varphi=[\underline\varphi_n,\overline\varphi_n]\) les cylindres
positifs et emboîtés produits par les deux lecteurs HMT\@. L'opération sur les
intervalles
\[
 \mathcal Q(I,J)
 :=
 \left[
 \frac{\inf I}{(\sup J)^2},
 \frac{\sup I}{(\inf J)^2}
 \right]
\]
est le plus petit intervalle qui contient \(x/y^2\) pour tout
\((x,y)\in I\times J\). Elle préserve l'emboîtement : si les cylindres d'entrée sont
raffinés, \(\mathcal Q(I_n^\pi,I_n^\varphi)\) est également raffiné. Comme leurs
diamètres tendent vers zéro et que la limite de \(I_n^\varphi\) est positive,
\[
 \bigcap_n\mathcal Q(I_n^\pi,I_n^\varphi)
 =
 \left\{\frac{\pi_{\rm HMT}}{\varphi_{\rm HMT}^2}\right\}.
\]
Avec les douze blocs en base mille déjà publiés, l'intervalle quotient impose
trente-cinq décimales :
\[
 \frac{\pi_{\rm HMT}}{\varphi_{\rm HMT}^2}
 =
 1.19998161486432666111577775331680692\ldots.
\]
Cette construction conserve la provenance de chaque chiffre : le numérateur
provient du cylindre de clôture et le dénominateur de l'autoéchelle.

\paragraph{Séparation algébrique entre lecture aréale, lecture volumétrique et leurs codomaines respectifs}
Le rapport complet ne peut provenir de la seule algèbre rationnelle de
\(F_{\rm av}\). En effet, ses fonctions rationnelles spectrales appartiennent à
\(\mathbb Q(\sqrt5)\) ; on y trouve \(\varphi^{-2}\), mais pas le nombre
transcendant \(\pi\). Par conséquent, toute dérivation qui prétendrait
obtenir \(\pi/\varphi^2\) uniquement à partir de la matrice \(2\times2\) serait
mal typée. La donnée transcendante doit provenir du lecteur coinductif de
clôture, ou d'un cocycle de frontière équivalent démontré
indépendamment. Ce contrôle négatif fixe la répartition des tâches entre
la correspondance finie et le lecteur de nombre.

\subsubsection{Involution d'échange du quotient}
\label{p12-83-c16-s2:sec:espejo-pi-phi-electron-hmt}

La correspondance admet deux orientations exactes d'une même loi. Si
\[
 \mathscr R(x,y)=\frac{x}{y^2},
 \qquad
 \mathsf J(x,y)=(y,x),
\]
alors \(\mathsf J^2=I\) et
\[
 \mathscr R(\pi,\varphi)=\frac{\pi}{\varphi^2},
 \qquad
 (\mathscr R\circ\mathsf J)(\pi,\varphi)
 =\frac{\varphi}{\pi^2}.
\]
On vérifie en outre les identités
\[
 \mathscr R(\pi,\varphi)\mathscr R(\varphi,\pi)
 =\frac1{\pi\varphi},
 \qquad
 \frac{\mathscr R(\pi,\varphi)}
 {\mathscr R(\varphi,\pi)}
 =\left(\frac{\pi}{\varphi}\right)^3.
\]
En écrivant
\(\mathcal B_{\rm av}^{(\pi)}=\mathscr R(\pi,\varphi)\), on obtient
également la reparamétrisation exacte
\[
 \frac{\varphi}{\pi^2}
 =
 \frac1{\varphi^3
 \bigl(\mathcal B_{\rm av}^{(\pi)}\bigr)^2},
 \qquad
 e^{\varphi/\pi^2}
 =
 \exp\!\left[
 \frac1{\varphi^3
 \bigl(\mathcal B_{\rm av}^{(\pi)}\bigr)^2}
 \right].
\]
Ainsi, le changement d'orientation entre clôture et croissance transporte le
rapport de frontière vers sa lecture échangée :
\[
 \frac{\pi}{\varphi^2}
 \xleftrightarrow{\ \mathsf J\ }
 \frac{\varphi}{\pi^2}.
\]
La première orientation lit la clôture aréale marquée sur la mémoire
volumétrique quadratique ; la seconde lit la croissance intérieure sur la
clôture quadratique. Cette correspondance est un théorème du lecteur HMT. La
mécanique dimensionnelle applique ensuite la seconde orientation dans la
construction électronique, sans faire de l'électron un antécédent de
\(\pi\), de \(\varphi\) ou de la correspondance.

\subsubsection{Principe d'Ambivalence des lecteurs aréal et volumétrique}
\label{p12-83-c16-s2:sec:ambivalencia-grav-inercial-pi-phi2}

Le Principe d'Ambivalence affirme que la lecture aréale et la lecture
volumétrique sont deux fonctionnelles coordonnées, mais non interchangeables, d'un
même état avec mémoire. Leur coïncidence sur la diagonale est
\[
 \mathcal L_{\rm ar}^{\rm diag}(x)
 =\mathcal L_{\rm vol}^{\rm diag}(x).
\]
L'égalité diagonale n'efface pas l'ambivalence en dehors de celle-ci.

Sur la frontière aréale--volumétrique, les deux lecteurs se factorisent à travers
une même amplitude positive de mémoire \(\mathcal M(x)\) :
\[
 \mathcal L_{\rm vol}(x)=\mu_{\rm vol}\mathcal M(x),
 \qquad
 \mathcal L_{\rm ar}(x)=
 \pi_{\rm HMT}\mu_{\rm ar}\mathcal M(x).
\]
Alors le théorème de Parry et la marque de retour donnent
\[
 \boxed{\displaystyle
 \frac{\mathcal L_{\rm ar}(x)}
 {\mathcal L_{\rm vol}(x)}
 =\frac{\pi_{\rm HMT}}{\varphi^2}.}
\]
Avec la normalisation volumétrique, la différence relative des lecteurs est
\[
 \frac{\mathcal L_{\rm ar}-\mathcal L_{\rm vol}}
 {\mathcal L_{\rm vol}}
 =
 \frac{\pi_{\rm HMT}}{\varphi^2}-1
 =
 0.1999816148643266611\ldots.
\]
La lecture physique de cette correspondance est effectuée plus tard, dans le chapitre de
Cartan--Holst, après la construction de l'holonomie, de la cotétrade, de la connexion, de la torsion
et de la courbure. Ici, le résultat est entièrement HMT : un rapport exact entre
deux lecteurs de la même mémoire de feuillet.

Il convient de les distinguer de l'extracteur enrichi utilisé ensuite en
mécanique dimensionnelle. Là, une histoire complète conserve le comptage d'action
\(n_A\), les douze événements orientés \(e_0,\ldots,e_{11}\) et la cochaîne
torsionnelle au niveau du pas ; on en obtient
\[
 s_C=\sum_{i=0}^{5}(e_i+e_{i+6}),
 \qquad
 W_\pm=6n_A\pm s_C.
\]
Les canaux déjà engendrés par HMT satisfont alors l'identité exacte
\[
 \exp\!\left(n_AA_{\rm rad}+\frac{s_C}{6}C^\ast_{\rm rad}\right)
 =\left(q_+^{-W_+}q_-^{-W_-}\right)^{1/12}.
\]
Par conséquent, le passage chemin--caractère angulaire n'est pas absent : il s'effectue sur
le registre dodécaphasique enrichi. Ce que ce calcul de phase prouve à lui
seul, ce sont \(\Omega_{\rm sw},S,G\) ; il ne les remplace pas et ne permet pas de reconstruire le
registre complet après l'avoir projeté.

La mesure uniforme sur les \(104\,976\) présentations descend vers des atomes
de masses \(1/729\), \(1/243\) et \(1/54\), mais pas vers \(1/468\). Dans la
microfibre quintuple de \(\pi\), les deux probabilités sont
\[
 \tau_{468}(\mathscr R_\pi)=\frac5{468},
 \qquad
 F_\#\mu_{\rm seed}(\mathscr R_\pi)=\frac7{729}.
\]
Le relèvement moyenné comme le relèvement de phase réalisent la première : ils ne
transportent pas la mesure uniforme sur les germes, mais attribuent la même masse
totale à chaque classe observable. Le relèvement de phase conserve en outre
l'immobilité microscopique dans les pas neutres.

Pour relier le relèvement de phase aux histoires enrichies de
profondeur six, il faut un noyau \(\mathcal K_6\) sur
\(\mathcal H_6^{\rm cat}\times C_9\). Celui-ci descend vers
\(\widehat{\mathcal K}_{\rm ph}\) à travers
\(\rho_6\times\mathrm{id}_{C_9}\) si et seulement s'il satisfait, à chaque phase, la
condition d'agrégabilité forte
\[
 \sum_{\rho_6(g)=y}\mathcal K_6((h,r),(g,r+1))
 =L_{\tau(r)}(\rho_6(h),y),
\]
indépendamment de \(h\) dans chaque fibre de \(\rho_6\). Dans ce cas,
la composition avec \(F\) se projette sur \(\mathcal K_{\rm ph}\). La construction
d'une famille \((\mathcal K_m)_{m\geq6}\) qui satisfait cette identité et
commute avec toutes les troncatures est l'étape restante vers un
transfert de la limite inverse.


\section{Correspondance aréale–volumétrique et rapport $\pi /\varphi ^2$}
\label{p12-83-c16-s3:sec:bisagra-hojas-pi-phi2}

La section précédente a construit l'involution dodécaphasique, son relèvement à des
sections persistantes, le ruban de Möbius relatif à une holonomie négative
et la carte biangulaire \((A,C^*)\leftrightarrow(q_+,q_-)\). Il reste maintenant à
exprimer précisément comment se comptent les retours des feuillets aréal et
volumétrique et pourquoi le nombre d'or apparaît au carré.

\subsection{Classes, supports et feuillets sont des objets différents}

Les classes nonadiques de mouvement sont
\[
 C_1=\{1,4,7\},\qquad C_2=\{2,5,8\},\qquad C_0=\{3,6,9\}.
\]
Les supports utilisés par l'observable angulaire sont
\[
 \mathcal D_{\rm act}=\{1,3,5,7\},\qquad
 \mathcal D_{\rm evt}=\{2,4,6,8\},\qquad
 \mathcal D_9=\{9\}.
\]
Leur croisement exact est
\[
 B_{\rm na}=\bigl(|C_i\cap D_j|\bigr)_{i,j}
 =\begin{pmatrix}2&1&0\\1&2&0\\1&1&1\end{pmatrix}.
\]

\begin{theorem}[Incidence triadique]
\label{p12-83-c16-s3:thm:incidencia-triadica-corta}
On a les identités
\[
 \det B_{\rm na}=3,
 \qquad
 \operatorname{SNF}(B_{\rm na})=\operatorname{diag}(1,1,3),
\]
\[
 \operatorname{im}B_{\rm na}
 =\{(u,v,w)\in\mathbb Z^3:u+v\equiv0\pmod3\}.
\]
L'échange simultané
\(C_1\leftrightarrow C_2\),
\(\mathcal D_{\rm act}\leftrightarrow\mathcal D_{\rm evt}\) laisse fixe
l'incidence.
\end{theorem}

\Needspace{5\baselineskip}
L'observable du registre
\[
 j_2(d)=
 \begin{cases}
 -1,&d\in\mathcal D_{\rm act},\\
 +1,&d\in\mathcal D_{\rm evt},\\
 0,&d\in\mathcal D_9
 \end{cases}
\]
se compose avec le lecteur géométrique actif :
\[
 -1\mapsto\mathscr H_{\rm ar},\qquad
 +1\mapsto\mathscr H_{\rm vol},\qquad
 0\mapsto\mathscr H_{\rm tor}.
\]
Par conséquent,
\[
 \mathfrak g_\triangle(\mathcal D_{\rm act})=\mathscr H_{\rm ar},\qquad
 \mathfrak g_\triangle(\mathcal D_{\rm evt})=\mathscr H_{\rm vol},\qquad
 \mathfrak g_\triangle(\mathcal D_9)=\mathscr H_{\rm tor}.
\]
La matrice prouve le croisement entre classes et supports. Le lecteur
\(\mathfrak g_\triangle\) déclare le passage des supports aux feuillets. Omettre cette
deuxième application reviendrait à identifier des cardinaux d'objets distincts.

\subsection{Ouverture spectrale du ruban au moyen de la phase angulaire conjuguée}

Dans la carte active,
\[
 q_+=e^{-A_{\rm rad}-C^*_{\rm rad}},\qquad
 q_-=e^{-A_{\rm rad}+C^*_{\rm rad}},
\]
et l'inversion exacte est
\[
 A_{\rm rad}=-\frac12\log(q_+q_-),
 \qquad
 C^*_{\rm rad}=\frac12\log\frac{q_-}{q_+}.
\]
Ainsi, \(A\) est le mode commun et \(C^*\) le mode différentiel. Pour chaque
hauteur de tour,
\[
 q_-^n-q_+^n
 =2e^{-nA_{\rm rad}}\sinh(nC^*_{\rm rad}).
\]
L'holonomie négative construite dans la
sous-section~\ref{sec:persistencia-ruta-mobius-familias} fixe le recollement
de Möbius ; la dernière identité démontre que \(C^*\) mesure l'ouverture
orientée des deux feuillets conjugués. Dans la représentation physique, ces deux
orientations se lisent comme particule et antiparticule.

Dans le lecteur de bord, un tour de la base est paramétré par \(2\pi\) et
se termine dans l'orientation opposée. Le lacet carré, paramétré par
\(4\pi\), rétablit l'orientation. C'est pourquoi, relativement au relèvement et au
lecteur géométrique déclarés, \(2\pi\) est la clôture aréale projetée et
\(4\pi\) la clôture volumétrique du revêtement orienté. Cette construction ne
confond pas ce double retour avec une représentation spinorielle ; la
représentation, lorsqu'elle est introduite, est un foncteur ultérieur.

\subsection{Descente de TPK, enveloppe des chemins et comptage des lacets}

Le calendrier TPK et l'observable angulaire \(j_2\) sont des applications distinctes. Le
premier met à jour le mode arithmétique
\(m\in\{\Sigma,\Pi\}\) ; le second classe les supports de chiffres. La
charnière utilise le lecteur HMT typé
\[
 \eta_{\rm av}(\Sigma)=\mathscr H_{\rm ar},
 \qquad
 \eta_{\rm av}(\Pi)=\mathscr H_{\rm vol}.
\]
Ainsi, \(\Sigma\) est le mode additif ou de bord et \(\Pi\) le mode
multiplicatif ou d'intérieur. Cette déclaration n'identifie pas
\(\tau\) à \(j_2\).

Dans le bloc primitif de phase
\[
 \tau_3=(+1,-1,0),
\]
TPK met à jour le mode par
\[
 +1:\ m\mapsto\Sigma,\qquad
 -1:\ m\mapsto\Pi,\qquad
 0:\ m\mapsto m.
\]
L'orbite cyclique est \((\Sigma,\Pi,\Pi)\). Définissons sa matrice
d'incidence des arêtes par
\[
 N_3(i,j)=
 \#\{r\in\mathbb Z/3\mathbb Z:(m_r,m_{r+1})=(i,j)\}.
\]

\begin{theorem}[Descente nécessaire des quatre règles]
\label{p12-83-c16-s3:thm:descenso-necesario-fav-corta}
Dans l’ordre \((\Sigma,\Pi)\),
\[
 N_3=
 \begin{pmatrix}0&1\\1&1\end{pmatrix}.
\]
De plus,
\[
 N_9=3N_3,\qquad
 N_{27}=9N_3,\qquad
 N_{108}=36N_3.
\]
Par conséquent, après avoir éliminé le facteur commun de répétition et appliqué
\(\eta_{\rm av}\), le quotient des modes de TPK induit nécessairement
\[
 \boxed{
 F_{\rm av}=S+P_{\rm vol}
 =
 \begin{pmatrix}0&1\\1&1\end{pmatrix}
 }
\]
où
\[
 S=\begin{pmatrix}0&1\\1&0\end{pmatrix},
 \qquad
 P_{\rm vol}=\begin{pmatrix}0&0\\0&1\end{pmatrix}.
\]
Les deux transitions croisées, l'unique autorétention volumétrique,
l'absence d'autorétention aréale et les multiplicités primitives unitaires
sont des conséquences de TPK, non des prémisses ajoutées.
\end{theorem}

\begin{proof}
Les trois couples consécutifs du cycle sont
\[
 \Sigma\longrightarrow\Pi,\qquad
 \Pi\longrightarrow\Pi,\qquad
 \Pi\longrightarrow\Sigma.
\]
Chacun apparaît une fois et \(\Sigma\to\Sigma\) n'apparaît pas. Cela détermine
les quatre entrées de \(N_3\). Les calendriers de longueurs
\(9,27,108\) contiennent respectivement \(3,9,36\) blocs primitifs ;
d'où les trois multiples. La réduction par le plus grand commun diviseur des
entrées non nulles n'élimine que la répétition du bloc. Le lecteur
\(\eta_{\rm av}\) change le type des sommets, non leur incidence.
\end{proof}

Le théorème précédent est une descente exacte d'arêtes. Il ne transforme pas l'oubli
de la phase en un quotient de Markov fort : une même émission visible peut
être mise à jour par \(I,E_+\) ou \(E_\times\), selon la phase cachée. De
ce fait découle une distinction nécessaire entre trois mesures :
\[
 \tau_{468},\qquad
 \nu_{\rm cyc},\qquad
 \mu_{\rm Parry}.
\]
La première est uniforme sur les \(468\) émissions et stationnaire pour le
transfert de phase. La seconde est la fréquence de l'orbite périodique,
\[
 \nu_{\rm cyc}(\mathscr H_{\rm ar})=\frac13,\qquad
 \nu_{\rm cyc}(\mathscr H_{\rm vol})=\frac23.
\]
Aucune d'elles ne peut produire directement le quotient irrationnel
\(\varphi^{-2}\).

L'objet qui le produit est déterminé par une propriété universelle.
Soit \(X_{\rm av}\) le plus petit sous-décalage de type fini à mémoire un qui
contient les mots obtenus en oubliant la phase et conserve tous leurs
couples consécutifs. Tout système à un pas contenant ces mots doit
admettre
\[
 {\rm ar}\to{\rm vol},\qquad
 {\rm vol}\to{\rm vol},\qquad
 {\rm vol}\to{\rm ar}.
\]
Le plus petit exclut l'unique couple absent, \({\rm ar}\to{\rm ar}\). Par
conséquent, \(X_{\rm av}\) a pour matrice d'adjacence \(F_{\rm av}\).
C'est la complétion symbolique minimale à mémoire un du quotient TPK ; ce
n'est pas une nouvelle dynamique microscopique.

\begin{theorem}[Loi de Fibonacci et mesure de Parry]
\label{p12-83-c16-s3:thm:ley-fibonacci-retornos-corta}
L'opérateur \(F_{\rm av}\) est primitif et
\[
 F_{\rm av}^{\,n}
 =\begin{pmatrix}F_{n-1}&F_n\\F_n&F_{n+1}\end{pmatrix}
 \qquad(n\geq1).
\]
Le sous-décalage \(X_{\rm av}\) a pour entropie topologique
\(\log\varphi\) et possède une unique mesure invariante d'entropie maximale. Pour
cette mesure,
\[
 \frac{(F_{\rm av}^{\,n})_{\rm ar,ar}}
 {(F_{\rm av}^{\,n})_{\rm vol,vol}}
 =\frac{F_{n-1}}{F_{n+1}}
 \longrightarrow\varphi^{-2},
 \qquad
 \boxed{\frac{\mu_{\rm ar}}{\mu_{\rm vol}}=\varphi^{-2}.}
\]
\end{theorem}

\begin{proof}
Le carré de \(F_{\rm av}\) est strictement positif. La formule des
puissances s'obtient par récurrence au moyen de
\(F_{n+1}=F_n+F_{n-1}\). Les entrées diagonales comptent les lacets fermés.
La racine de Perron est \(\varphi\) et un vecteur de Perron positif est
\((1,\varphi)\). Comme la matrice est symétrique, les poids de sommet de
Parry sont proportionnels au produit des vecteurs gauche et droit,
c'est-à-dire à \((1,\varphi^2)\).
\end{proof}

Le facteur \(\varphi^{-1}\) est le rapport des amplitudes de Perron d'une
traversée ; \(\varphi^{-2}\) est le rapport des poids d'un retour complet.
Les deux exposants sont ainsi ordonnés sans introduire deux constantes
indépendantes.

\subsection{Origine spectrale de l'exposant quadratique du retour}

L'apparition du carré admet une seconde démonstration, purement
spectrale, qui ne dépend pas de l'interprétation d'une moyenne pondérée. Le polynôme
caractéristique de \(F_{\rm av}\) est
\[
 \chi_{F_{\rm av}}(t)=t^2-t-1,
\]
de sorte que ses deux valeurs propres sont
\[
 \lambda_+=\varphi,\qquad \lambda_-=-\varphi^{-1}.
\]
Le signe de \(\lambda_-\) enregistre l'inversion d'orientation d'une
traversée du feuillet stable. En relevant l'action à la mémoire quadratique
\(\operatorname{Sym}^2(\mathbb R^2)\), les valeurs propres sont les produits
par paires
\[
 \lambda_+^2=\varphi^2,\qquad
 \lambda_+\lambda_-=-1,\qquad
 \lambda_-^2=\varphi^{-2}.
\]
Dans la base
\((e_{\rm ar}^2,e_{\rm ar}\odot e_{\rm vol},e_{\rm vol}^2)\), en écrivant
par colonnes les coordonnées des images des trois vecteurs de base, la
matrice est
\[
 \operatorname{Sym}^2(F_{\rm av})=
 \begin{pmatrix}
 0&0&1\\
 0&1&2\\
 1&1&1
 \end{pmatrix},
\]
et satisfait
\[
 \chi_{\operatorname{Sym}^2(F_{\rm av})}(t)
 =t^3-2t^2-2t+1
 =(t+1)(t^2-3t+1).
\]
Par conséquent, le mode stable change d'orientation avec
\(-\varphi^{-1}\) en une traversée et devient le poids positif
\(\varphi^{-2}\) lorsqu'il est lu comme mémoire d'aire ou comme retour
quadratique. C'est la raison structurelle de l'exposant deux.

\subsection{Clôture de bord et quotient \texorpdfstring{$\pi/\varphi^2$}{pi/phi2}}

Le lecteur coinductif de \(\pi\) précède cette construction. Nous définissons le
ruban de clôture par composition, sans utiliser \(\pi\) pour sélectionner le
transporteur :
\[
 \mathcal B_{\rm av}^{(\pi)}(X_\infty)
 :=\Pi_{\rm HMT}(X_\infty)
 \frac{\mu_{\rm ar}}{\mu_{\rm vol}}.
\]

\paragraph{Retour marqué \(\Theta_n\) et convergence du rapport aréal–volumétrique}
Soient les projecteurs de feuillet
\[
 P_{\rm ar}=\begin{pmatrix}1&0\\0&0\end{pmatrix},
 \qquad
 P_{\rm vol}=\begin{pmatrix}0&0\\0&1\end{pmatrix}.
\]
Le marqueur de frontière est défini par
\[
 D_\pi:=\Pi_{\rm HMT}(X_\infty)P_{\rm ar}+P_{\rm vol}.
\]
C'est l'unique opérateur positif diagonal par rapport à la décomposition en
feuillets qui attribue un poids unitaire au feuillet volumétrique et la clôture
\(\Pi_{\rm HMT}(X_\infty)\) au feuillet aréal. Son caractère canonique est donc
relatif à trois clauses explicites : diagonalité de feuillet, normalisation
volumétrique et marquage aréal par le lecteur de clôture déjà construit.

Si \(e_{\rm ar},e_{\rm vol}\) sont les vecteurs de feuillet, la fonctionnelle de
retour marqué de profondeur \(n\) est
\[
 \Theta_n(X_\infty):=
 \frac{\langle e_{\rm ar},
 D_\pi F_{\rm av}^{\,n}e_{\rm ar}\rangle}
 {\langle e_{\rm vol},
 F_{\rm av}^{\,n}e_{\rm vol}\rangle}.
\]
La loi des puissances précédente donne, sans approximation,
\[
 \boxed{
 \Theta_n(X_\infty)
 =\Pi_{\rm HMT}(X_\infty)\frac{F_{n-1}}{F_{n+1}}.}
\]
La clôture \(\pi\) est appliquée une seule fois, lors du marquage du retour de bord ; elle n'est pas
réintroduite à chaque itération de \(F_{\rm av}\).

\begin{corollary}[Rapport holographique aréal--volumétrique]
\label{p12-83-c16-s3:cor:pi-phi2-corta}
Relativement au quotient des modes TPK, à son enveloppe minimale à mémoire
un et à la clôture HMT de \(\pi\),
\[
 \boxed{\mathcal B_{\rm av}^{(\pi)}
 =\lim_{n\to\infty}\Theta_n
 =\frac{\pi_{\rm HMT}}{\varphi^2}.}
\]
\end{corollary}

\paragraph{Périodicité native de \(108\) phases et clôture du calendrier unitaire}
À la profondeur dodécaphasique \(108\), les deux comptages diagonaux sont
\[
 (F_{\rm av}^{108})_{\rm ar,ar}
 =F_{107}=10284720757613717413913,
\]
\[
 (F_{\rm av}^{108})_{\rm vol,vol}
 =F_{109}=26925748508234281076009.
\]
Par conséquent,
\[
 \frac{F_{107}}{F_{109}}
 =0.381966011250105151795413165634361882\ldots,
\]
et l'erreur par rapport à \(\varphi^{-2}\) est
\[
 \left|\frac{F_{107}}{F_{109}}-\varphi^{-2}\right|
 =6.1684979916614407936\ldots\times10^{-46}.
\]
Après le marquage de bord, l'erreur par rapport à
\(\pi/\varphi^2\) est
\[
 \left|\pi\frac{F_{107}}{F_{109}}-\frac{\pi}{\varphi^2}\right|
 =1.9378907974286976068\ldots\times10^{-45}.
\]
Ainsi, \(108\) n'est pas utilisé pour ajuster le rapport : c'est une profondeur native du
calendrier qui approche déjà le retour limite avec plus de quarante-quatre
chiffres fractionnaires stables.

\paragraph{Lecteur d'intervalles coinductif : emboîtement des cylindres et unicité de la limite}
Pour des intervalles positifs
\[
 I=[a,b],\qquad J=[c,d],\qquad 0<c\leq d,
\]
nous définissons
\[
 \mathcal Q(I,J):=
 \left[\frac{a}{d^2},\frac{b}{c^2}\right].
\]
La fonction \((x,y)\mapsto x/y^2\) est croissante en \(x\) et décroissante en
\(y>0\). C'est pourquoi \(\mathcal Q(I,J)\) est exactement le plus petit intervalle
qui contient tous les quotients \(x/y^2\) avec
\((x,y)\in I\times J\) ; ce n'est pas une borne choisie par arrondi.

Soient \(I_n^\pi\) et \(I_n^\varphi\) les cylindres positifs emboîtés
engendrés par HMT\@. Alors
\[
 \mathcal Q(I_{n+1}^\pi,I_{n+1}^\varphi)
 \subseteq
 \mathcal Q(I_n^\pi,I_n^\varphi),
\]
et, comme
\[
 \bigcap_n I_n^\pi=\{\pi_{\rm HMT}\},
 \qquad
 \bigcap_n I_n^\varphi=\{\varphi_{\rm HMT}\},
\]
on obtient
\[
 \bigcap_n\mathcal Q(I_n^\pi,I_n^\varphi)
 =\left\{\frac{\pi_{\rm HMT}}{\varphi_{\rm HMT}^2}\right\}.
\]
Chaque bloc de base \(1000\) est imposé dès que l'intervalle quotient
entre entièrement dans un unique cylindre triadique de la profondeur
correspondante.

En particulier, les cylindres positifs de douze triades pour \(\pi\) et
\(\varphi\) produisent l'intervalle quotient
\[
 \left[
 \frac{p^-}{(f^+)^2},
 \frac{p^+}{(f^-)^2}
 \right].
\]
L'évaluation certifie trente-cinq chiffres fractionnaires :
\[
 \frac{\pi}{\varphi^2}
 =1.19998161486432666111577775331680692\ldots,
\]
et onze blocs complets :
\[
 199|981|614|864|326|661|115|777|753|316|806.
\]

\paragraph{Théorème de séparation algébrique.}
La matrice \(F_{\rm av}\), ses puissances et son algèbre spectrale rationnelle
vivent dans une extension algébrique de \(\mathbb Q\), contenue dans
\(\mathbb Q(\sqrt5)\). Par conséquent, aucune composition utilisant uniquement
des opérations algébriques rationnelles sur \(F_{\rm av}\) ne peut produire le
facteur transcendant \(\pi/\varphi^2\). L'opérateur fini engendre
\(\varphi^{-2}\) ; le lecteur coinductif de frontière engendre
\(\pi_{\rm HMT}\) ; le retour marqué compose les deux sans confondre leurs
provenances. Ce contrôle négatif empêche de présenter \(F_{\rm av}\) comme
générateur isolé de \(\pi\).

\subsection{Involution des lecteurs et section électronique}
\label{p12-83-c16-s3:sec:espejo-pi-phi-electron}

La relation qui vient d'être obtenue apparaît auprès de l'électron d'une manière plus
précise qu'une répétition littérale. L'invariant électronique actif de la
sous-section~\ref{sec:persistencia-ruta-mobius-familias} contient
\[
 \mathfrak e_0=
 \frac{\sqrt3}{4}
 \left(e^{\varphi/\pi^2}-22\alpha_{\rm HMT}^{3}\right)
 (1+15\Delta_4).
\]
Par conséquent, l'exposant électronique est \(\varphi/\pi^2\), tandis que la
charnière des feuillets produit \(\pi/\varphi^2\). Les deux scalaires ne sont
ni égaux ni inverses l'un de l'autre. Ils forment les deux orientations d'une même loi
bivariée. En effet, définissons
\[
 \mathscr R(x,y):=\frac{x}{y^2},
 \qquad
 \mathsf J(x,y):=(y,x).
\]
Alors \(\mathsf J^2=I\) et
\[
 \boxed{
 \mathscr R(\pi,\varphi)=\frac{\pi}{\varphi^2},
 \qquad
 (\mathscr R\circ\mathsf J)(\pi,\varphi)
 =\frac{\varphi}{\pi^2}.}
\]
L'échange des deux lecteurs conduit exactement le retour marqué
aréal--volumétrique au germe exponentiel de l'électron. En outre,
\[
 \mathscr R(\pi,\varphi)\mathscr R(\varphi,\pi)
 =\frac1{\pi\varphi},
 \qquad
 \frac{\mathscr R(\pi,\varphi)}
 {\mathscr R(\varphi,\pi)}
 =\left(\frac{\pi}{\varphi}\right)^3.
\]
Si nous notons
\(\mathcal B_{\rm av}^{(\pi)}=\pi/\varphi^2\) le rapport marqué des
feuillets, la même identité peut s'écrire sans introduire de constante
supplémentaire :
\[
 \boxed{
 \frac{\varphi}{\pi^2}
 =
 \frac{1}
 {\varphi^3\bigl(\mathcal B_{\rm av}^{(\pi)}\bigr)^2},}
 \qquad
 e^{\varphi/\pi^2}
 =
 \exp\!\left[
 \frac{1}
 {\varphi^3\bigl(\mathcal B_{\rm av}^{(\pi)}\bigr)^2}
 \right].
\]
Numériquement,
\[
 \begin{aligned}
 \frac{\varphi}{\pi^2}
 &=0.163941118913672383599752755745666982\ldots,\\
 e^{\varphi/\pi^2}
 &=1.17814494259630678081160095680888910\ldots.
 \end{aligned}
\]

Cette identité formalise le contenu bidirectionnel sans modifier la force de
l'équation électronique. L'orientation \((\pi,\varphi)\) lit la clôture de bord
par mémoire quadratique ; l'orientation échangée
\((\varphi,\pi)\) lit la croissance intérieure par clôture quadratique et entre
exponentiée dans \(\mathfrak e_0\). Le théorème démontre que les deux
expressions appartiennent à la même orbite d'échange. Les coefficients
\(22\) et \(15\) possèdent un générateur indépendant : le
\cref{p12-83-c15-s2:thm:selector-persistencia-electron-rev7} les obtient, sans consulter la
masse électronique, comme
\[
 -22=-\bigl|\mathcal K_e\setminus\iota(\mathscr R_\pi)\bigr|,
 \qquad
 +15=\bigl|\mathscr R_\pi\times\mathbb F_3\bigr|.
\]
Leur statut est \textsc{démontré avec structure de départ explicite}.
L'involution actuelle transporte cette construction vers l'orientation électronique et
n'utilise pas l'électron pour obtenir \(F_{\rm av}\), \(\pi\) ou \(\varphi\). La
dernière forme est un changement exact de coordonnées du germe électronique, non
une dérivation alternative de \(\mathfrak e_0\).

\subsection{Identité pentadique de Rogers--Ramanujan et reconnaissance de l'invariant de Perron}

La fraction continue de Rogers--Ramanujan fournit une réalisation
analytique du même invariant. Avec
\(q_A=e^{-A_{\rm rad}}\), \(q_C=e^{-C^*_{\rm rad}}\), on obtient
\[
 R(q_A)R(q_C)-\varphi^{-2}\simeq-1.01864\times10^{-27},
\]
et avec les canaux de la décomposition scindée,
\[
 R(q_+)R(q_-)-\varphi^{-2}\simeq-1.19458\times10^{-21}.
\]
Les différences ne sont pas nulles. L'affirmation exacte est la limite
\[
 q_A^{(m)},q_C^{(m)}\to1^-
 \quad\Longrightarrow\quad
 R(q_A^{(m)})R(q_C^{(m)})\to\varphi^{-2}.
\]
Rogers--Ramanujan reconnaît ainsi l'invariant de Perron dans le lecteur
pentadique ; elle ne sélectionne pas à elle seule les coordonnées angulaires \(A\), \(C^*\) ni
\(F_{\rm av}\). L'ancien noyau fini \(90/120\) conserve un résidu
approché \(7.81\times10^{-9}\) et n'est pas utilisé comme identité exacte.

\subsection{Involution des lecteurs inertiel et gravitationnel et reconstruction bidirectionnelle de la grandeur}

Le principe d'Ambivalence distingue le lecteur d'ouverture
\(q_A\), le lecteur de rétention \(q_V\) et la mémoire transportée. La
charnière précédente permet de formuler exactement l'intuition physique, sans
la confondre avec le théorème combinatoire. Soit \(\mathcal M(x)\) un caractère
commun de mémoire et définissons, à l'interface aréale--volumétrique,
\[
 \mathcal I_{\rm av}(x)
 :=\mu_{\rm vol}\mathcal M(x),
 \qquad
 \mathcal G_{\rm av}(x)
 :=\pi_{\rm HMT}\mu_{\rm ar}\mathcal M(x).
\]
Alors, pourvu que \(\mathcal M(x)\ne0\),
\[
 \boxed{
 \frac{\mathcal G_{\rm av}(x)}{\mathcal I_{\rm av}(x)}
 =\frac{\pi_{\rm HMT}}{\varphi^2},}
 \qquad
 \frac{\mathcal G_{\rm av}-\mathcal I_{\rm av}}
 {\mathcal I_{\rm av}}
 =\frac{\pi_{\rm HMT}}{\varphi^2}-1.
\]
La première formule exprime le rapport orienté entre les lecteurs
gravitationnel et inertiel ; la seconde, leur défaut relatif. Le L gnomonique
porte la carte euclidienne locale sur laquelle les deux lecteurs descendent vers une
normalisation commune,
\[
 \frac{\mathcal G_{L}}{\mathcal I_{L}}=1.
\]

La lecture physique est stratifiée. Dans la carte euclidienne du L, la
connexion torsionnelle locale se projette sur une normalisation unique et le
principe d'équivalence se réalise comme quotient un. Dans l'état
enrichi antérieur à cette projection, l'ouverture aréale et la rétention
volumétrique restent séparées : le Principe d'Ambivalence conserve leurs
deux fonctionnelles et fixe le quotient interne \(\pi/\varphi^2\). L'équivalence
locale et l'ambivalence de frontière sont ainsi deux lectures successives du
même antécédent, avec des domaines et des applications déclarés.

La réalisation physique adopte la factorisation de l'observable inertiel par la
rétention volumétrique et celle de l'observable gravitationnel par l'ouverture aréale
avec clôture de bord \(\pi\) ; les deux partagent le caractère \(\mathcal M\). La
mathématique interne détermine la correspondance et son quotient, et le foncteur de
mécanique dimensionnelle publie ses cartes observables. Une rupture de l'une quelconque
de ces factorisations constitue le falsificateur explicite de cette réalisation.

\paragraph{Frontière cosmologique.}
L'hypothèse selon laquelle la matière noire ou l'énergie noire seraient des lectures effectives
de croisements non torsionnels est conservée uniquement comme
hypothèse de prolongement physique. Pour la transformer en proposition
testable,
il faut construire au moins une application des interfaces de feuillet vers une
géométrie cosmologique, les observables représentant la densité et la pression et
une comparaison prospective avec des données sans utiliser celles-ci pour fixer
l'application. La correspondance démontrée fournit le candidat structurel ; elle ne
constitue pas encore à elle seule une explication observationnelle du secteur sombre.

\paragraph{Domaine de la charnière aréale–volumétrique et conditions nécessaires de validité}
Le ruban, le rapport aréal--volumétrique et la lecture bidirectionnelle appartiennent à
l'architecture antérieure. La matrice d'incidence, la descente du calendrier,
l'enveloppe minimale, la loi des lacets, la mesure de Parry, l'action
quadratique et le retour marqué fournissent sa formalisation explicite. Le
lecteur d'intervalles et les contrôles de profondeur \(108\) certifient le
calcul. Les incidences \(N_3,N_9,N_{27},N_{108}\), les identités
matricielles et les cylindres sont des identités internes exactes ; le théorème de
Möbius et le rapport
\(\pi/\varphi^2\) sont
\textsc{démontrés avec structure de départ explicite} ; les lectures
particule--antiparticule et inertielle--gravitationnelle constituent des interfaces
physiques ultérieures ; l'extension au secteur sombre reste une
hypothèse physique non démontrée. Le couple relié par l'involution
\(\mathscr R(\pi,\varphi)\leftrightarrow
\mathscr R(\varphi,\pi)\) formalise exactement une architecture antérieure. Les
coefficients \((-22,+15)\) sont démontrés, par le
\cref{p12-83-c15-s2:thm:selector-persistencia-electron-rev7}, avec leur structure de départ
explicite.

Une holonomie positive donne un cylindre. Une transition aréale--aréale,
l'absence d'une transition croisée ou un comptage TPK différent de
\(N_9=3F_{\rm av}\) et \(N_{108}=36F_{\rm av}\) falsifient la descente.
L'identification de \(\mu_{\rm Parry}\) à \(\tau_{468}\) ou à
\(\nu_{\rm cyc}\) est réfutée par leurs valeurs distinctes. Un spectre de
\(\operatorname{Sym}^2(F_{\rm av})\) différent de
\(\{\varphi^2,-1,\varphi^{-2}\}\), un comptage \(108\) différent de
\((F_{107},F_{109})\), un intervalle quotient non emboîté ou l'introduction
itérative de \(\pi\) dans \(F_{\rm av}\) falsifient les formalisations
nouvelles. Une identification physique universelle du quotient sans application des
feuillets aux observables reste hors du théorème. Le delta des quatre règles
est clos ; il ne rouvre pas \(\pi\), \(\varphi\), \(C^*\) ni le raccord
extension--chemin--registre--caractère.
